%% ---------------------------------------------------------
\section{Site-dependent chains at fixed length}
\label{sec:non_ti_setup}
%% ---------------------------------------------------------

The algebraic proof of~\cite{Molnar2018NormalPEPS} uses periodic square chains
whose tensors vary from site to site. The proof of
\cite[Theorem~3]{PerezGarcia2007Matrix} starts instead from a varying-bond
open chain. We record both fixed-length constructions and distinguish their
boundary conventions.

\begin{definition}[Periodic non-TI chain data]\label{def:non_ti_chain}
    \lean{MPSChainTensor, MPSChainTensor.coeff, MPSChainTensor.SameState, MPSChainTensor.IsInjective, MPSChainTensor.GaugeEquiv}
    \leanok
    \uses{def:mps_tensor, def:injective}
    Fix a chain length $n$.
    A \emph{non-translation-invariant periodic MPS chain} consists of sitewise
    tensors $A_k=\{A_k^i\}_{i=0}^{d-1}$, where $k=0,\ldots,n-1$ and each
    $A_k^i\in\MN{D}$.
    For a basis configuration
    $\sigma=(\sigma_0,\ldots,\sigma_{n-1})$, its coefficient is
    \begin{align}
        c_{A,\sigma}
         & :=\tr\!\left(A_0^{\sigma_0}A_1^{\sigma_1}\cdots A_{n-1}^{\sigma_{n-1}}\right).
        \notag
    \end{align}
    Two such chains generate the \emph{same chain state} if these coefficients
    agree for every $\sigma$.
    The chain is \emph{injective} if each local tensor $A_k$ is injective in
    the usual single-site sense.
    Finally, two chains are \emph{cyclically gauge equivalent} if there are
    invertible bond matrices $Z_0,\ldots,Z_{n-1}$ such that, for every site
    $k$ and physical index $i$,
    \begin{align}
        B_k^i
         & =Z_kA_k^iZ_{k+1}^{-1},
        \notag
    \end{align}
    with indices understood modulo $n$, so that $Z_n=Z_0$. Locally, this
    gauge relation is represented by
    \begin{center}
        % B_k^i = Z_k A_k^i Z_{k+1}^{-1}: the black nodes are the
        % site tensors, the red capsules are the two neighboring bond
        % gauges, and the upper leg carries the free physical index i.
        % Both sides have boundary signature (1,1,1,0), ordered as
        % (virtual-west, virtual-east, physical-up, physical-down).
        \tenkzkernel
        \begin{tenkz}[cols=1, west=open, east=open]
            \tn[ports={90:physical:$i$}]{B_k}
        \end{tenkz}
        $ = $
        \begin{tenkz}[west=open, east=open]
            \tn[skin=ring]{Z_k} & \tn[ports={90:physical:$i$}]{A_k} & \tn[skin=ring]{Z_{k+1}^{-1}}
        \end{tenkz}
    \end{center}
\end{definition}

\begin{definition}[Varying-bond open chain]
    \label{def:varying_bond_obc_chain}
    \lean{OBCChainTensor}
    \leanok
    Fix a chain length $N$ and dimensions $D_0,\ldots,D_N\leq D$ with
    $D_0=D_N=1$. A varying-bond open chain consists of matrices
    $A_k^i\in\mathbb C^{D_k\times D_{k+1}}$ for
    $k=0,\ldots,N-1$ and $i=0,\ldots,d-1$. Intermediate dimensions may
    vanish. The number $D$ is a common upper bound and need not be the least
    such bound.

    This is the open-boundary representation of
    \cite[Section~II.B]{PerezGarcia2007Matrix}. The source takes
    $D=\max_k D_k$; allowing an arbitrary common upper bound does not change
    the rectangular chain.
\end{definition}

\begin{definition}[Open-chain coefficient]
    \label{def:varying_bond_obc_coeff}
    \lean{OBCChainTensor.coeff}
    \leanok
    \uses{def:varying_bond_obc_chain}
    For a physical configuration
    $\sigma=(\sigma_0,\ldots,\sigma_{N-1})$, define
    \begin{align}
        c_A(\sigma)
         & :=\sum_{\alpha_0=0}^{D_0-1}\cdots
        \sum_{\alpha_N=0}^{D_N-1}
        \prod_{k=0}^{N-1}
        (A_k^{\sigma_k})_{\alpha_k,\alpha_{k+1}}.
        \label{eq:varying_bond_obc_coeff}
    \end{align}
    Since $D_0=D_N=1$, this is the scalar matrix product in
    \cite[Section~II.B]{PerezGarcia2007Matrix}. If an intermediate dimension
    vanishes, the virtual-path set is empty. For $N=0$, there is one empty
    path and its empty product is $1$.
\end{definition}

\begin{definition}[Cyclic translation of an open-chain configuration]
    \label{def:varying_bond_obc_translation}
    \lean{OBCChainTensor.cyclicTranslateConfig}
    \leanok
    Let $N\geq1$. For $s\in\mathbb Z/N\mathbb Z$, the cyclic translation of
    a physical configuration is
    \begin{align}
        (T_s\sigma)_k & =\sigma_{k+s}.
        \label{eq:varying_bond_obc_translation}
    \end{align}
    Only the physical configuration is translated. The dimensions
    $D_0,\ldots,D_N$, the site matrices, and the distinguished open cut remain
    fixed. At $N=0$, there is no translation parameter.
\end{definition}

\begin{definition}[Translation invariance at fixed length]
    \label{def:varying_bond_obc_ti}
    \lean{OBCChainTensor.IsTranslationInvariantState}
    \leanok
    \uses{def:varying_bond_obc_coeff, def:varying_bond_obc_translation}
    For $N\geq1$, a varying-bond open chain represents a
    translation-invariant state at length $N$ if, for every
    $s\in\mathbb Z/N\mathbb Z$ and every physical configuration $\sigma$,
    \begin{align}
        c_A(T_s\sigma) & =c_A(\sigma).
        \label{eq:varying_bond_obc_ti}
    \end{align}
    This condition compares coefficients at one fixed length. It does not
    rotate the bond dimensions or identify open cuts. At $N=0$, the condition
    is vacuous.
\end{definition}

\begin{theorem}[Empty open-chain coefficient]
    \label{thm:varying_bond_obc_coeff_zero}
    \lean{OBCChainTensor.coeff_zero}
    \leanok
    \uses{def:varying_bond_obc_coeff}
    At $N=0$, the open-chain coefficient is $1$.
\end{theorem}

\begin{proof}
    \leanok
    \uses{def:varying_bond_obc_coeff}
    The left endpoint space is one-dimensional, so there is one virtual path.
    Its product contains no factors and therefore equals $1$.
\end{proof}

\begin{theorem}[Vanishing from a zero virtual dimension]
    \label{thm:varying_bond_obc_zero_bond}
    \lean{OBCChainTensor.coeff_eq_zero_of_bondDim_eq_zero}
    \leanok
    \uses{def:varying_bond_obc_coeff}
    If $D_k=0$ for some $k\in\{0,\ldots,N\}$, then
    $c_A(\sigma)=0$ for every physical configuration $\sigma$.
\end{theorem}

\begin{proof}
    \leanok
    \uses{def:varying_bond_obc_coeff}
    No virtual path can have a coordinate in the empty $k$th bond space, so
    the sum in~(\ref{eq:varying_bond_obc_coeff}) is empty.
\end{proof}

\begin{definition}[Square padding of a varying-bond open chain]
    \label{def:varying_bond_obc_zero_pad}
    \lean{OBCChainTensor.zeroPad}
    \leanok
    \uses{def:varying_bond_obc_chain, def:non_ti_chain}
    For each site $k$, embed $A_k^i$ into the upper-left
    $D_k\times D_{k+1}$ corner of a $D\times D$ matrix
    $\widetilde A_k^i$ and set every other entry to zero. The matrices
    $\widetilde A_k^i$ form a square site-dependent closed chain.

    This square padding is an auxiliary algebraic factorization.
    The proof of \cite[Theorem~3]{PerezGarcia2007Matrix} does not first replace
    the rectangular open chain by this square chain. It places the original
    rectangular matrices directly into the compatible off-diagonal blocks of
    one $ND\times ND$ cyclic matrix.
\end{definition}

\begin{theorem}[Empty padded-chain coefficient]
    \label{thm:varying_bond_obc_zero_pad_zero}
    \lean{OBCChainTensor.coeff_zeroPad_zero}
    \leanok
    \uses{def:varying_bond_obc_zero_pad, def:non_ti_chain}
    At $N=0$, the coefficient of the padded square chain is $D$.
\end{theorem}

\begin{proof}
    \leanok
    \uses{def:varying_bond_obc_zero_pad, def:non_ti_chain}
    The empty matrix product is the $D\times D$ identity, whose trace is $D$.
\end{proof}

\begin{theorem}[Zero-extension of a dependent finite sum]
    \label{thm:pi_fin_sum_zero_extension}
    \lean{Fintype.sum_piFin_castLE_extend_zero}
    \leanok
    Let $I$ be finite, let $r_i\leq s_i$ be nonnegative integers for
    $i\in I$, and let $f$ take values in an additive commutative monoid. Then
    the sum of $f$ over $\prod_{i\in I}\{0,\ldots,r_i-1\}$ equals its sum over
    $\prod_{i\in I}\{0,\ldots,s_i-1\}$ after extending $f$ by zero outside
    the smaller Cartesian product.
\end{theorem}

\begin{proof}
    \leanok
    The coordinatewise inclusions identify the smaller Cartesian product with
    the subset of the larger product on which every coordinate is below its
    original bound. Split the larger sum into this subset and its complement.
\end{proof}

\begin{theorem}[Cyclic path expansion for a square chain]
    \label{thm:chain_coeff_sum_cyclic}
    \lean{MPSChainTensor.coeff_eq_sum_cyclic}
    \leanok
    \uses{def:non_ti_chain}
    Let $N\geq1$ and let $B_0,\ldots,B_{N-1}$ be a square site-dependent
    chain of bond dimension $D$. For every physical configuration $\sigma$,
    \begin{align}
        c_{B,\sigma}
         & =\sum_{g\colon\mathbb Z/N\mathbb Z\to\{0,\ldots,D-1\}}
        \prod_{k\in\mathbb Z/N\mathbb Z}
        (B_k^{\sigma_k})_{g(k),g(k+1)}.
        \label{eq:chain_coeff_sum_cyclic}
    \end{align}
\end{theorem}

\begin{proof}
    \leanok
    \uses{lem:eval_word_of_fn_prod, thm:trace_eval_word_eq_sum_cyclic}
    Encode the site and physical indices as one enlarged physical index and
    apply the cyclic trace expansion for a closed matrix product.
\end{proof}

\begin{theorem}[Positive-length coefficient preservation]
    \label{thm:varying_bond_obc_zero_pad_coeff}
    \lean{OBCChainTensor.coeff_zeroPad}
    \leanok
    \uses{def:varying_bond_obc_coeff, def:varying_bond_obc_zero_pad}
    Let $N\geq1$. For every physical configuration $\sigma$,
    \begin{align}
        \tr\!\left(\widetilde A_0^{\sigma_0}\cdots
        \widetilde A_{N-1}^{\sigma_{N-1}}\right)
         & =c_A(\sigma).
        \label{eq:varying_bond_obc_zero_pad_coeff}
    \end{align}
    No positivity assumption is imposed on the intermediate dimensions. The
    endpoint conditions already imply $D\geq1$ whenever such a chain exists.
    The equality concerns only this fixed length.
\end{theorem}

\begin{proof}
    \leanok
    \uses{thm:chain_coeff_sum_cyclic, thm:pi_fin_sum_zero_extension, def:varying_bond_obc_coeff, def:varying_bond_obc_zero_pad}
    Expand the left-hand side by~(\ref{eq:chain_coeff_sum_cyclic}). Extend the
    dependent sum in~(\ref{eq:varying_bond_obc_coeff}) to the common dimension
    $D$ by Theorem~\ref{thm:pi_fin_sum_zero_extension}. Every extended path
    outside the rectangular bond ranges has a zero matrix entry. For the
    remaining paths, the one-dimensional endpoint spaces identify the last
    virtual coordinate with the first, and the products agree term by term.
\end{proof}

\begin{theorem}[Padded coefficient with a zero virtual dimension]
    \label{thm:varying_bond_obc_zero_pad_zero_bond}
    \lean{OBCChainTensor.coeff_zeroPad_eq_zero_of_bondDim_eq_zero}
    \leanok
    \uses{def:non_ti_chain, def:varying_bond_obc_zero_pad}
    Let $N\geq1$. If $D_k=0$ for some
    $k\in\{0,\ldots,N\}$, then the padded square-chain coefficient vanishes
    for every physical configuration.
\end{theorem}

\begin{proof}
    \leanok
    \uses{thm:varying_bond_obc_zero_pad_coeff, thm:varying_bond_obc_zero_bond}
    Apply positive-length coefficient preservation and then the vanishing of
    the open-chain contraction.
\end{proof}
\begin{definition}[Cyclic block tensor at fixed length]
    \label{def:chain_cyclic_block_tensor}
    \lean{MPSChainTensor.cyclicRotation, MPSChainTensor.cyclicBlockTensor}
    \leanok
    \uses{def:non_ti_chain}
    Let $N\geq 1$ and let $A_0,\ldots,A_{N-1}$ be a periodic chain of
    $D\times D$ matrices. For $k\in\mathbb Z/N\mathbb Z$, write
    $(A^{(k)})_j=A_{j+k}$ for the cyclic rotation by $k$ sites. Let $P$ be
    the permutation matrix on the block coordinate satisfying
    $P(e_a\otimes e_k)=e_a\otimes e_{k-1}$. For each physical index $i$, set
    \begin{align}
        B^i
         & =N^{-1/N}\operatorname{diag}
        \left(A_0^i,\ldots,A_{N-1}^i\right)P
        \in\MN{ND}.
        \label{eq:chain_cyclic_block_tensor}
    \end{align}
    Thus the $(k,k+1)$ block of $B^i$ is $N^{-1/N}A_k^i$, with block
    indices taken modulo $N$, and every other block vanishes. The tensor $B$
    depends on the fixed length $N$.
\end{definition}

\begin{theorem}[Fixed-length cyclic block average]
    \label{thm:chain_cyclic_block_average}
    \lean{MPSChainTensor.mpv_cyclicBlockTensor_eq_average}
    \leanok
    \uses{def:chain_cyclic_block_tensor, def:non_ti_chain}
    For every basis configuration $\sigma\in\{0,\ldots,d-1\}^N$, the
    length-$N$ coefficient of the tensor in
    Definition~\ref{def:chain_cyclic_block_tensor} is
    \begin{align}
        V^{(N)}(B)_\sigma
         & =\frac{1}{N}\sum_{k=0}^{N-1}c_{A^{(k)},\sigma}.
        \label{eq:chain_cyclic_block_average}
    \end{align}
    Reindexing the finite sum and cyclically permuting factors under the trace
    identifies this with the physical-configuration rotation in
    \cite[Theorem~3]{PerezGarcia2007Matrix}. This identity concerns only the
    fixed length $N$.
\end{theorem}

\begin{proof}
    \leanok
    \uses{def:chain_cyclic_block_tensor, def:non_ti_chain, lem:eval_word_smul}
    Write $D_i=\operatorname{diag}(A_0^i,\ldots,A_{N-1}^i)$ and omit the
    scalar normalization temporarily. Commuting the cyclic permutation through
    a block diagonal gives
    \begin{align}
        P\operatorname{diag}(M_0,\ldots,M_{N-1})
         & =\operatorname{diag}(M_1,\ldots,M_0)P.
        \label{eq:chain_cyclic_block_commute}
    \end{align}
    Repeated use of~(\ref{eq:chain_cyclic_block_commute}), followed by
    $P^N=I$, yields
    \begin{align}
        D_{\sigma_0}P\cdots D_{\sigma_{N-1}}P
         & =\operatorname{diag}\left(
        \prod_{j=0}^{N-1}A_{j+k}^{\sigma_j}
        \right)_{k=0}^{N-1}.
        \label{eq:chain_cyclic_block_product}
    \end{align}
    Taking the trace of~(\ref{eq:chain_cyclic_block_product}) sums the traces
    of its diagonal blocks. Finally,
    $(N^{-1/N})^N=N^{-1}$ gives
    (\ref{eq:chain_cyclic_block_average}).
\end{proof}

% Scope restriction (supplied OBC representation): PGVWC07 Theorem 3 also
% obtains this input from purity via its OBC representation theorem. See
% docs/paper-gaps/pgvwc07_site_independent_scope.tex.
\begin{theorem}[Conditional site-independent matrices at fixed length]
    \label{thm:pgvwc07_site_independent_matrices}
    \lean{OBCChainTensor.pgvwc07_site_independent_matrices,
        OBCChainTensor.exists_pgvwc07_site_independent_tensor}
    \leanok
    \uses{def:varying_bond_obc_ti, def:varying_bond_obc_zero_pad, def:chain_cyclic_block_tensor}
    This is the cyclic-block construction in
    \cite[Theorem~3]{PerezGarcia2007Matrix}, conditional on its open-chain
    input. Let $N\geq1$. Suppose a translation-invariant state on the $N$-site
    ring has a varying-bond open-chain representation whose virtual dimensions are
    at most $D$. Then it has a site-independent trace-MPS representation of
    bond dimension $ND$ with the same length-$N$ coefficients.
    The constructed tensor depends on $N$, and the equality is asserted only
    at this ring length.
\end{theorem}

\begin{proof}
    \leanok
    \uses{thm:chain_coeff_sum_cyclic, thm:varying_bond_obc_zero_pad_coeff, thm:chain_cyclic_block_average}
    The cyclic rotation of the padded square chain satisfies
    \begin{align}
        c_{(\widetilde A)^{(k)},\sigma}
         & =c_A(\tau_{-k}\sigma),
        \label{eq:pgvwc07_zero_pad_rotation}
    \end{align}
    where $(\tau_s\sigma)_j=\sigma_{j+s}$. To prove
    (\ref{eq:pgvwc07_zero_pad_rotation}), expand both square-chain traces as
    cyclic virtual-path sums and reindex the virtual path by the same rotation.
    The inverse shift appears when the rotated matrix product is restored to
    the distinguished open cut. Translation invariance makes every term in
    the cyclic average equal to $c_A(\sigma)$, so the factor $1/N$ cancels the
    sum of $N$ identical terms.
\end{proof}

\begin{lemma}[Cyclic gauge equivalence is an equivalence relation]
    \label{lem:chain_gauge_equiv_equivalence}
    \lean{MPSChainTensor.GaugeEquiv.refl}
    \lean{MPSChainTensor.GaugeEquiv.symm}
    \lean{MPSChainTensor.GaugeEquiv.trans}
    \leanok
    \uses{def:non_ti_chain}
    On fixed-length periodic chains, cyclic gauge equivalence is reflexive,
    symmetric, and transitive.
\end{lemma}

\begin{proof}
    \leanok
    Reflexivity uses the identity gauge on every bond.
    Symmetry inverts each bond gauge, and transitivity multiplies the gauges
    pointwise along the chain.
\end{proof}

%% ---------------------------------------------------------
\section{Quadratic reconstruction of the translation-invariant canonical form}%
\label{sec:quadratic_reconstruction}
%% ---------------------------------------------------------

A translation-invariant state whose successive reduced density operators have
rank at most $D^2$ can be stored as an open-boundary matrix product state with
$D^2\times D^2$ matrices. The reconstruction theorem of
\cite[Theorem~8]{PerezGarcia2007Matrix} recovers a translation-invariant
tensor of bond dimension $D$ from a window of $D^4$ consecutive matrices of
the open-boundary canonical form by solving a system of quadratic equations
whose size does not depend on the chain length. This section records the
system and the two auxiliary results that drive the uniqueness half of the
reconstruction argument, both stated in the source immediately before the
proof of the uniqueness theorem \cite[Theorem~7]{PerezGarcia2007Matrix}.

\begin{definition}[Solution of the quadratic reconstruction system]
    \label{def:pgvwc07_quadratic_system}
    \lean{MPSTensor.QuadraticReconstructionSolution}
    \leanok
    Let a window of matrix families on the doubled bond space be given,
    $A^{[j]}_i\in\MN{D^2}$ for $i\in\Fin{d}$ and $0\le j\le m-1$. A solution
    of the quadratic reconstruction system over this window consists of
    matrices $Y_0,\ldots,Y_{m-1}$ and $Z_0,\ldots,Z_{m-1}$ in $\MN{D^2}$
    together with matrices $B_i\in\MN{D}$ satisfying
    \begin{align}
        B_i\otimes\Id           & = Y_j A^{[j]}_i Z_j &  & \forall\, i,j,
        \label{eq:quadratic_system_gauge}                                      \\
        Y_{j+1}Z_j              & = \Id               &  & \forall\, j\le m-2,
        \label{eq:quadratic_system_inverse}                                    \\
        \sum_i B_iB_i^{\dagger} & = \Id.
        \label{eq:quadratic_system_normalization}
    \end{align}
    In the system~(S) of \cite{PerezGarcia2007Matrix} the window consists of
    the $m=D^4$ matrices of the unique open-boundary canonical form at the
    sites $L_0+1,\ldots,L_0+D^4$; the unknown $Y_j$ above is the source
    unknown $Y_{L_0+1+j}$, while $Z_j$ is the source unknown
    $Z_{L_0+2+j}$ attached to the right of that site.
\end{definition}

\begin{lemma}[Invertibility of the paired gauges]
    \label{lem:pgvwc07_quadratic_system_units}
    \lean{MPSTensor.QuadraticReconstructionSolution.isUnit_leftGauge,
        MPSTensor.QuadraticReconstructionSolution.isUnit_rightGauge}
    \leanok
    \uses{def:pgvwc07_quadratic_system}
    In any solution of the quadratic reconstruction system, every matrix
    $Y_{j+1}$ and every matrix $Z_j$ appearing
    in~(\ref{eq:quadratic_system_inverse}) is invertible.
\end{lemma}

\begin{proof}
    \leanok
    A square matrix with a one-sided inverse is invertible.
\end{proof}

\begin{lemma}[Nonzero combination along a shifted chain]
    \label{lem:pgvwc07_chain_combination}
    \lean{pgvwc07ChainCoeff, pgvwc07ChainCoeff_succ,
        pgvwc07_chainCoeff_combination_spec}
    \leanok
    Let $T,S$ be linear maps defined on the same vector spaces, and let
    $Y_0,\ldots,Y_n$ be vectors such that $T(Y_k)=S(Y_{k+1})$ for every
    $0\le k\le n-1$, the vectors $Y_0,\ldots,Y_{n-1}$ are linearly
    independent, and
    \begin{align}
        Y_n & = \sum_{k=0}^{n-1}\lambda_k Y_k.
        \label{eq:chain_combination_dependence}
    \end{align}
    Let $x\neq0$ be a solution of the polynomial equation
    \begin{align}
        \sum_{k=0}^{n-1}\lambda_k x^{n-k} & = 1,
        \label{eq:chain_combination_polynomial}
    \end{align}
    and define $\mu_j=\sum_{k=0}^{j}\lambda_k x^{j+1-k}$. Then the
    combination $Y=\sum_{k=0}^{n-1}\mu_k Y_k$ is nonzero and satisfies
    $T(Y)=\tfrac1x S(Y)$. This is the first auxiliary result before the proof
    of \cite[Theorem~7]{PerezGarcia2007Matrix}, written with indices starting
    at zero.
\end{lemma}

\begin{proof}
    \leanok
    The coefficients satisfy $\mu_0=\lambda_0 x$ and the recurrence
    $\mu_{j+1}=x(\lambda_{j+1}+\mu_j)$, and
    (\ref{eq:chain_combination_polynomial}) states $\mu_{n-1}=1$. Applying
    $T$ termwise and using the chain relation,
    \begin{align}
        T(Y)
         & =\sum_{k=0}^{n-1}\mu_k S(Y_{k+1})
        =S\Bigl(\sum_{k=0}^{n-2}\mu_k Y_{k+1}+Y_n\Bigr)
        =S\Bigl(\sum_{j=0}^{n-1}(\lambda_j+\mu_{j-1})Y_j\Bigr),
        \label{eq:chain_combination_reindex}
    \end{align}
    with the convention $\mu_{-1}=0$, where the last equality substitutes
    (\ref{eq:chain_combination_dependence}). By the recurrence each
    coefficient equals $\mu_j/x$, so the right-hand side is
    $\tfrac1x S(Y)$. Since $\mu_{n-1}=1$ and $Y_0,\ldots,Y_{n-1}$ are
    linearly independent, $Y\neq0$.
\end{proof}

\begin{lemma}[Amplified intertwiners factor through the multiplicity]
    \label{lem:pgvwc07_kronecker_intertwiner}
    \lean{Matrix.kroneckerSlice,
        Matrix.kronecker_one_intertwines_iff_kroneckerSlice_intertwines,
        Matrix.exists_nonzero_intertwiner_of_kronecker_one_intertwines}
    \leanok
    Let $B,C\in\MN{n}$ and let $W$ act on the doubled space. Then
    \begin{align}
        W(C\otimes\Id) & = (B\otimes\Id)W
        \label{eq:kronecker_intertwiner_equation}
    \end{align}
    holds if and only if every multiplicity slice $X_{ab}$ of $W$, with
    entries $(X_{ab})_{pq}=W_{(p,a),(q,b)}$, satisfies $X_{ab}C=BX_{ab}$.
    Consequently a nonzero simultaneous solution $W$ of the amplified
    equations for a family of pairs $(B_i,C_i)$ yields a nonzero matrix $X$
    with $XC_i=B_iX$ for every $i$. This realizes the second auxiliary result
    before the proof of \cite[Theorem~7]{PerezGarcia2007Matrix}, which
    phrases the solution space of~(\ref{eq:kronecker_intertwiner_equation})
    as $\mathcal{S}\otimes M_n$ for $\mathcal{S}$ the solution space of
    $XC=BX$.
\end{lemma}

\begin{proof}
    \leanok
    Evaluating both sides of~(\ref{eq:kronecker_intertwiner_equation}) at a
    pair of doubled indices reduces the identity factor to a Kronecker delta,
    which collapses the double sum to the slice products $X_{ab}C$ and
    $BX_{ab}$. For the consequence, a nonzero $W$ has a nonzero slice.
\end{proof}

\begin{lemma}[Chain of intertwiners between two solutions]
    \label{lem:pgvwc07_solution_chain}
    \lean{MPSTensor.QuadraticReconstructionSolution.chainIntertwiner,
        MPSTensor.QuadraticReconstructionSolution.isUnit_chainIntertwiner,
        MPSTensor.QuadraticReconstructionSolution.chainIntertwiner_mul_kronecker}
    \leanok
    \uses{def:pgvwc07_quadratic_system, lem:pgvwc07_quadratic_system_units}
    Let two solutions of the quadratic reconstruction system over the same
    window be given, with row unknowns $Y_j$, $Y'_j$, column unknowns $Z_j$,
    $Z'_j$, and tensors $B_i$, $C_i$. For every window site $j$ carrying an
    inversion equation, the matrix
    \begin{align}
        W_{j+1} & = Y'_{j+1}Z_j = Y'_{j+1}Y_{j+1}^{-1}
        \label{eq:solution_chain_intertwiner}
    \end{align}
    is invertible, and whenever the neighbouring site also carries an
    inversion equation,
    \begin{align}
        W_{j+1}(B_i\otimes\Id) & = (C_i\otimes\Id)W_{j+2} &  & \forall\, i.
        \label{eq:solution_chain_relation}
    \end{align}
    Thus a window of $m$ sites supplies exactly $m-1$ internal
    intertwiners and $m-2$ relations between consecutive internal
    intertwiners. This is the internal part of the chain produced at the start
    of the proof of \cite[Theorem~8]{PerezGarcia2007Matrix}. It does not by
    itself provide the linear dependence required by
    Lemma~\ref{lem:pgvwc07_chain_combination}.
\end{lemma}

\begin{proof}
    \leanok
    Both solutions express $B_i\otimes\Id$ and $C_i\otimes\Id$ through the
    same window matrix at site $j+1$; substituting one expression into the
    other and cancelling the paired gauges with
    (\ref{eq:quadratic_system_inverse}) gives
    (\ref{eq:solution_chain_relation}). Invertibility follows from
    Lemma~\ref{lem:pgvwc07_quadratic_system_units}.
\end{proof}

\begin{theorem}[Obtaining the translation-invariant canonical form]
    \label{thm:pgvwc07_quadratic_reconstruction}
    \notready
    \uses{def:pgvwc07_quadratic_system}
    Consider a translation-invariant state on an $N$-site ring with unique
    open-boundary canonical form such that the rank of each reduced density
    operator of successive spins is bounded by $D^2$, and form the quadratic
    system of Definition~\ref{def:pgvwc07_quadratic_system} over the window of
    $D^4$ open-boundary canonical matrices starting after the injectivity
    onset. Write $A=(A_i)_i$ for its canonical translation-invariant tensor.
    \begin{enumerate}
        \item If the state admits a translation-invariant matrix product
              representation satisfying condition C1, then the system has a
              solution.
        \item Let $B=(B_i)_i$ be the tensor of a solution and suppose in
              addition that its length-$N$ trace-MPS representation is the
              given state. Then there are a unitary $U$ and a scalar
              $\xi\in\C$ such that $|\xi|=1$, $\xi^N=1$, and
              \begin{align}
                  A_i & =\xi UB_iU^{\dagger}
                  \qquad\text{for every }i.
                  \label{eq:pgvwc07_quadratic_reconstruction_phase_conjugacy}
              \end{align}
              Consequently the explicitly rephased tensor
              $\widehat B_i:=\xi B_i$ still represents the same length-$N$
              state and is exactly unitarily conjugate to $A$.
    \end{enumerate}
    No assertion is made here that every algebraic solution represents the
    prescribed state. The unrestricted representation claim and the exact
    conjugacy claimed in the source have a residual finite-ring phase defect
    \cite{gap:pgvwc07_quadratic_reconstruction_phase}. This is the corrected
    form of \cite[Theorem~8]{PerezGarcia2007Matrix}; its formalization is not
    yet complete.
\end{theorem}

\begin{proof}
    \uses{lem:pgvwc07_chain_combination, lem:pgvwc07_kronecker_intertwiner, lem:pgvwc07_quadratic_system_units, lem:pgvwc07_solution_chain}
    The canonical representation solves the system by the freedom in the
    open-boundary representation. For the converse, assume that the solution
    tensor $B$ represents the given length-$N$ state. Comparing its
    open-boundary representation with the canonical one supplies intertwiners
    at all $D^4+1$ cuts around the $D^4$-site window, including both endpoint
    cuts, together with the $D^4$ adjacent relations
    \cite{gap:pgvwc07_intertwiner_chain_off_by_one}. Linear dependence in
    $M_{D^2}(\C)$ then provides the hypothesis of
    Lemma~\ref{lem:pgvwc07_chain_combination}. The internal construction in
    Lemma~\ref{lem:pgvwc07_solution_chain} gives only $D^4-1$ intertwiners
    and $D^4-2$ relations, so it does not force this dependence for an
    arbitrary algebraic solution.

    Once the full cut chain is supplied, the combination lemma and the slice
    extraction of Lemma~\ref{lem:pgvwc07_kronecker_intertwiner} give a
    nonzero matrix $R$ and a scalar $x\neq0$ with
    $RB_i=\tfrac1x A_iR$. The two normalizations
    (\ref{eq:quadratic_system_normalization}), the invariance of the positive
    diagonal fixed point, and the uniqueness of the transfer-map fixed point
    under condition C1 give $|x|=1$ and
    $RR^{\dagger}=c\Id$ for some $c>0$. Thus $U=c^{-1/2}R$ is unitary and
    $A_i=xUB_iU^{\dagger}$ for every $i$. Since $A$ and $B$ represent the
    same nonzero length-$N$ state, their coefficients satisfy
    $V^{(N)}(A)=x^NV^{(N)}(B)=V^{(N)}(B)$, and hence $x^N=1$. The tensor
    $\widehat B_i=xB_i$ therefore represents the same state and satisfies
    $A_i=U\widehat B_iU^{\dagger}$.
\end{proof}

%% ---------------------------------------------------------
\section{Uniqueness of the translation-invariant canonical form at a fixed ring length}%
\label{sec:pgvwc07_ti_uniqueness}
%% ---------------------------------------------------------

The uniqueness theorem \cite[Theorem~7]{PerezGarcia2007Matrix} compares two
translation-invariant canonical representations of one state on a ring of
$N$ sites. Its printed proof rewrites both representations as open-boundary
representations on the doubled bond space, compares them with the
open-boundary canonical form, and obtains invertible intertwiners at the
cuts of a window of $D^4$ sites; the two auxiliary results of
Section~\ref{sec:quadratic_reconstruction} then produce a nonzero
intertwiner between the two local tensors. This section records the
cut-vector argument that supplies those intertwiners directly from condition
C1, the resulting fixed-length intertwiner, and the corrected uniqueness
theorem. The printed conclusion is false on a finite ring because of an
$N$th-root phase \cite{gap:pgvwc07_quadratic_reconstruction_phase}; the
corrected statement carries that phase explicitly
\cite{gap:pgvwc07_ti_uniqueness_scope}.

\begin{definition}[Cut vectors]
    \label{def:pgvwc07_cut_vectors}
    \lean{MPSTensor.leftCutVector, MPSTensor.rightCutVector,
        MPSTensor.coeff_append_eq_sum_cutVector,
        MPSTensor.leftCutVector_append_single}
    \leanok
    \uses{def:mps_tensor}
    Let $A=(A_i)_i$ be a tensor of bond dimension $D$ and let $L\ge0$. For
    $\alpha,\beta\in\Fin{D}$ the \emph{left cut vector} and the \emph{right
        cut vector} of length $L$ are the vectors
    \begin{align}
        \lvert\Phi_{\alpha,\beta}\rangle
         & =\sum_{i_1,\ldots,i_L}\langle\alpha|A_{i_1}\cdots A_{i_L}|\beta\rangle
        \,\lvert i_1\cdots i_L\rangle,
        \label{eq:pgvwc07_left_cut_vector}                                        \\
        \lvert\Psi_{\alpha,\beta}\rangle
         & =\sum_{i_1,\ldots,i_L}\langle\beta|A_{i_1}\cdots A_{i_L}|\alpha\rangle
        \,\lvert i_1\cdots i_L\rangle.
        \label{eq:pgvwc07_right_cut_vector}
    \end{align}
    For a word split into a left part of length $k$ and a right part of
    length $m$, the coefficient of the ring state factors across the cut as
    \begin{align}
        \operatorname{tr}\bigl(A_{i_1}\cdots A_{i_k}A_{j_1}\cdots A_{j_m}\bigr)
         & =\sum_{\alpha,\beta}\Phi_{\alpha,\beta}(i_1\cdots i_k)\,
        \Psi_{\alpha,\beta}(j_1\cdots j_m),
        \label{eq:pgvwc07_cut_resolution}
    \end{align}
    and extending the left word by one site multiplies the left cut vectors
    by the matrix of that site,
    $\Phi_{\alpha,\beta}(i_1\cdots i_k\,i)
        =\sum_\gamma\Phi_{\alpha,\gamma}(i_1\cdots i_k)\,(A_i)_{\gamma\beta}$.
    These are the vectors introduced in the proof of the proposition on
    condition C1 in \cite{PerezGarcia2007Matrix}.
\end{definition}

\begin{lemma}[Linear independence of the cut vectors]
    \label{lem:pgvwc07_cut_vectors_linearIndependent}
    \lean{MPSTensor.linearIndependent_leftCutVector,
        MPSTensor.linearIndependent_rightCutVector,
        MPSTensor.isNBlkInjective_of_le_of_unital}
    \leanok
    \uses{def:pgvwc07_cut_vectors, def:l_blk_injective}
    If $A$ is $L$-block injective, then the $D^2$ left cut vectors of length
    $L$ are linearly independent, and so are the $D^2$ right cut vectors of
    length $L$. If moreover $\sum_iA_iA_i^\dagger=\Id$, then $A$ is
    $L'$-block injective for every $L'\ge L$, so both statements hold at
    every length at least $L$.
\end{lemma}

\begin{proof}
    \leanok
    A vanishing combination $\sum_{\alpha,\beta}c_{\alpha,\beta}
        \lvert\Phi_{\alpha,\beta}\rangle=0$ says that the linear functional
    $M\mapsto\sum_{\alpha,\beta}c_{\alpha,\beta}M_{\alpha\beta}$ annihilates
    every product of $L$ matrices of $A$. These products span $M_D(\C)$, so
    the functional vanishes on the matrix units and $c=0$. The right cut
    vectors are treated in the same way with the transposed functional.
    For the propagation, every matrix $M$ equals
    $\sum_iA_i(A_i^\dagger M)$, and each $A_i^\dagger M$ is a combination of
    products of $L'$ matrices, so $M$ is a combination of products of $L'+1$
    matrices.
\end{proof}

\begin{lemma}[Intertwiners at the cuts]
    \label{lem:pgvwc07_cut_intertwiner}
    \lean{MPSTensor.exists_cutGauge_of_coeff_eq, MPSTensor.cutGauge_chain}
    \leanok
    \uses{def:pgvwc07_cut_vectors}
    Let $B$ and $C$ be tensors of bond dimension $D$ whose coefficients agree
    on every word of length $k+m$, and suppose that the left cut vectors of
    $B$ of length $k$ and the right cut vectors of $B$ of length $m$ are
    linearly independent. Then there is an invertible matrix $W$ on the
    doubled bond space with
    \begin{align}
        \Phi^{C}_{a}
         & =\sum_{b}W_{b,a}\,\Phi^{B}_{b}
        \qquad\text{for every index pair }a,
        \label{eq:pgvwc07_cut_intertwiner}
    \end{align}
    where the index pairs are written column first. If $W$ and $W'$ realize
    (\ref{eq:pgvwc07_cut_intertwiner}) at the cuts after $k$ and after $k+1$
    sites, and the left cut vectors of $B$ of length $k$ are linearly
    independent, then
    \begin{align}
        W(C_i\otimes\Id) & =(B_i\otimes\Id)W' \qquad\text{for every }i.
        \label{eq:pgvwc07_cut_chain_relation}
    \end{align}
    This is the square case of the freedom theorem for open-boundary
    representations \cite[Theorem~2]{PerezGarcia2007Matrix} invoked in the
    proof of \cite[Theorem~7]{PerezGarcia2007Matrix}, obtained here by a
    direct comparison of the two bipartite decompositions
    (\ref{eq:pgvwc07_cut_resolution}).
\end{lemma}

\begin{proof}
    \leanok
    \uses{lem:pgvwc07_cut_vectors_linearIndependent,
        thm:peps_existsUnique_virtualOperation_of_endPair}
    Both decompositions (\ref{eq:pgvwc07_cut_resolution}) of the same
    coefficients have $D^2$ terms. The dual functionals of the linearly
    independent right cut vectors of $B$ read off a matrix $G$ with
    $\Phi^{B}_{\mu}=\sum_\nu G_{\mu\nu}\Phi^{C}_{\nu}$; linear independence of
    the left cut vectors of $B$ shows that no nonzero row vector annihilates
    $G$, so $G$ is invertible, and $W=(G^{-1})^{\mathsf T}$ satisfies
    (\ref{eq:pgvwc07_cut_intertwiner}). For the chain relation, extend the
    left words by one site: the two expressions of $\Phi^{C}$ at length
    $k+1$ through (\ref{eq:pgvwc07_cut_intertwiner}) at the two cuts are
    combinations of the independent family $\Phi^{B}$ of length $k$, and
    comparing coefficients gives (\ref{eq:pgvwc07_cut_chain_relation})
    entrywise.
\end{proof}

\begin{theorem}[Fixed-length intertwiner]
    \label{thm:pgvwc07_fixed_length_intertwiner}
    \lean{MPSTensor.exists_intertwiner_of_fixedLength_mpv_eq,
        MPSTensor.exists_pgvwc07_chain_root}
    \leanok
    \uses{def:l_blk_injective, def:mpv}
    Let $B$ be a tensor of bond dimension $D\ge1$ with
    $\sum_iB_iB_i^\dagger=\Id$ that is $L_0$-block injective, let
    $N\ge2L_0+D^4$, and let $C$ be a tensor of bond dimension $D$ with the
    same length-$N$ coefficients as $B$. Then there are a nonzero matrix
    $R\in\MN{D}$ and a scalar $x\neq0$ such that
    \begin{align}
        RC_i & =\tfrac1x\,B_iR \qquad\text{for every }i.
        \label{eq:pgvwc07_fixed_length_intertwiner}
    \end{align}
    This is the intertwiner obtained in the proof of
    \cite[Theorem~7]{PerezGarcia2007Matrix} before the unitarity argument,
    with the intertwiner count corrected to the $D^4+1$ cuts of the window
    \cite{gap:pgvwc07_intertwiner_chain_off_by_one}.
\end{theorem}

\begin{proof}
    \leanok
    \uses{lem:pgvwc07_cut_vectors_linearIndependent,
        lem:pgvwc07_cut_intertwiner, lem:pgvwc07_chain_combination,
        lem:pgvwc07_kronecker_intertwiner}
    Consider the $D^4+1$ cuts after $L_0,\ldots,L_0+D^4$ sites; each leaves
    at least $L_0$ sites on both sides, so by
    Lemma~\ref{lem:pgvwc07_cut_vectors_linearIndependent} the cut vectors of
    $B$ are linearly independent there, and
    Lemma~\ref{lem:pgvwc07_cut_intertwiner} supplies invertible matrices
    $W_0,\ldots,W_{D^4}$ on the doubled bond space with
    $W_k(C_i\otimes\Id)=(B_i\otimes\Id)W_{k+1}$. Since $\MN{D^2}$ has
    dimension $D^4$, there is a first index $n\ge1$ such that
    $W_0,\ldots,W_{n-1}$ are linearly independent and
    $W_n=\sum_{k<n}\lambda_kW_k$; some $\lambda_k$ is nonzero because $W_n$
    is invertible, so the polynomial equation
    (\ref{eq:chain_combination_polynomial}) has positive degree and a nonzero
    root $x$. Lemma~\ref{lem:pgvwc07_chain_combination}, applied to right
    multiplication by $C_i\otimes\Id$ and left multiplication by
    $B_i\otimes\Id$, gives a nonzero combination $W$ of $W_0,\ldots,W_{n-1}$
    with $W(C_i\otimes\Id)=(\tfrac1xB_i\otimes\Id)W$ for every $i$, and
    Lemma~\ref{lem:pgvwc07_kronecker_intertwiner} extracts a nonzero
    multiplicity slice $R$ satisfying (\ref{eq:pgvwc07_fixed_length_intertwiner}).
\end{proof}

\begin{lemma}[Unitarity of the intertwiner]
    \label{lem:pgvwc07_intertwiner_unitary}
    \lean{MPSTensor.isUnit_of_intertwines_of_isNBlkInjective,
        MPSTensor.exists_unitaryConj_of_intertwines_of_isNBlkInjective,
        MPSTensor.mpv_eq_pow_mul_mpv_of_unitaryConj}
    \leanok
    \uses{def:l_blk_injective, def:mpv}
    Let $B$ and $C$ satisfy $\sum_iB_iB_i^\dagger=\Id$ and
    $\sum_iC_iC_i^\dagger=\Id$, let $B$ be $L_0$-block injective, and let
    $R\neq0$ and $x\neq0$ satisfy (\ref{eq:pgvwc07_fixed_length_intertwiner}).
    Then $R$ is invertible, $|x|=1$, and there is a unitary $U$ with
    \begin{align}
        B_i & =x\,UC_iU^\dagger \qquad\text{for every }i.
        \label{eq:pgvwc07_intertwiner_unitary_conjugacy}
    \end{align}
    Consequently the length-$N$ coefficients of $B$ and $C$ satisfy
    $V^{(N)}(B)=x^NV^{(N)}(C)$ for every $N$.
\end{lemma}

\begin{proof}
    \leanok
    \uses{lem:left_canonical_irreducible_same_phase_unitary_gauge}
    The range of $R$ is invariant under every $B_i$, hence under every
    product of $L_0$ matrices of $B$; block injectivity leaves no nontrivial
    invariant subspace, so the range is everything and $R$ is invertible.
    Block injectivity also makes the transfer maps of $B$ and of $C$
    irreducible, the latter because $C_i=\tfrac1xR^{-1}B_iR$. The
    conjugate-transposed families are therefore left-canonical irreducible
    tensors related by the gauge $(R^\dagger)^{-1}$ and the scalar
    $\overline{x}$, and
    Lemma~\ref{lem:left_canonical_irreducible_same_phase_unitary_gauge}
    replaces the gauge by a unitary and gives $|x|=1$. Taking adjoints
    returns (\ref{eq:pgvwc07_intertwiner_unitary_conjugacy}). This is the
    fixed-point argument of \cite[Theorem~7]{PerezGarcia2007Matrix}, which
    derives $|x|=1$ from the dual fixed point of $B$ and $RR^\dagger=\Id$
    from the single-block property; both follow from condition C1 through
    the Perron--Frobenius theory of irreducible maps. The coefficient
    relation follows from (\ref{eq:pgvwc07_intertwiner_unitary_conjugacy})
    because conjugation by $U$ does not change traces.
\end{proof}

\begin{definition}[Translation-invariant canonical representation]
    \label{def:pgvwc07_ti_canonical_representation}
    \lean{MPSTensor.PGVWC07CanonicalFormData, MPSTensor.IsPGVWC07CanonicalForm}
    \leanok
    \uses{def:mps_tensor}
    A tensor $A$ of bond dimension $D$ is a \emph{translation-invariant
        canonical representation} when its bond indices split into finitely
    many blocks of positive size, and there are weights $1\ge\lambda_j>0$ and
    tensors $A^j$ on the blocks with
    \begin{align}
        A_i & =\bigoplus_j\lambda_jA_i^j \qquad\text{for every }i,
        \label{eq:pgvwc07_ti_canonical_representation_block_form}
    \end{align}
    such that every block satisfies the three conditions of
    Theorem~\ref{thm:pgvwc07_ti_canonical_form}:
    $\sum_iA_i^j(A_i^j)^\dagger=\Id$, there is a diagonal positive definite
    $\Lambda^j$ with $\sum_i(A_i^j)^\dagger\Lambda^jA_i^j=\Lambda^j$, and
    the identity is the only fixed point of
    $X\mapsto\sum_iA_i^jX(A_i^j)^\dagger$ up to scalars. This is the block
    form of \cite[Theorem~4]{PerezGarcia2007Matrix}, which lists the blocks
    consecutively.
\end{definition}

\begin{lemma}[Canonical representations from the existence theorem]
    \label{lem:pgvwc07_ti_canonical_representation_of_blocks}
    \lean{MPSTensor.isPGVWC07CanonicalForm_toTensorFromBlocks}
    \leanok
    \uses{def:pgvwc07_ti_canonical_representation,
        thm:pgvwc07_ti_canonical_form}
    The weighted direct sum $\bigoplus_k\mu_kA_k$ of positive-dimensional
    blocks produced by Theorem~\ref{thm:pgvwc07_ti_canonical_form}, with
    positive weights $\mu_k\le1$, is a translation-invariant canonical
    representation with the blocks in consecutive order.
\end{lemma}

\begin{proof}
    \leanok
    The weights are positive reals of modulus at most one, and the block
    conditions are those of Definition~\ref{def:pgvwc07_ti_canonical_representation}.
\end{proof}

\begin{lemma}[A block-injective canonical representation has one block]
    \label{lem:pgvwc07_ti_canonical_representation_single_block}
    \lean{MPSTensor.PGVWC07CanonicalFormData.evalWord_eq,
        MPSTensor.PGVWC07CanonicalFormData.not_isNBlkInjective_of_ne,
        MPSTensor.PGVWC07CanonicalFormData.sum_mul_conjTranspose_eq,
        MPSTensor.PGVWC07CanonicalFormData.exists_weight_sum_mul_conjTranspose_eq_of_isNBlkInjective}
    \leanok
    \uses{def:pgvwc07_ti_canonical_representation, def:l_blk_injective}
    Let $A$ be a translation-invariant canonical representation with the
    block form (\ref{eq:pgvwc07_ti_canonical_representation_block_form}).
    Every product $A_{i_1}\cdots A_{i_L}$ is block diagonal, with
    $\lambda_j^LA^j_{i_1}\cdots A^j_{i_L}$ on block $j$, and
    \begin{align}
        \sum_iA_iA_i^\dagger & =\bigoplus_j\lambda_j^2\,\Id.
        \label{eq:pgvwc07_ti_canonical_representation_sum}
    \end{align}
    If $A$ has two distinct blocks, then $A$ is not $L$-block injective for
    any $L$. Consequently, if $D\ge1$ and $A$ is $L$-block injective for
    some $L$, then $A$ has a single block, and
    $\sum_iA_iA_i^\dagger=\lambda^2\,\Id$ with $0<\lambda\le1$. This is the
    first assertion of the proposition on condition C1 in
    \cite{PerezGarcia2007Matrix}.
\end{lemma}

\begin{proof}
    \leanok
    The block-diagonal form of the products follows from
    (\ref{eq:pgvwc07_ti_canonical_representation_block_form}) by induction
    on the length, and (\ref{eq:pgvwc07_ti_canonical_representation_sum})
    from the first block condition. If $\alpha$ lies in one block and
    $\beta$ in another, the entry $(\alpha,\beta)$ of every product
    vanishes while the matrix unit $\lvert\alpha\rangle\langle\beta\rvert$
    has entry one there, so the products of any length do not span
    $M_D(\C)$. When $D\ge1$ there is at least one block, and block
    injectivity leaves exactly one; then
    (\ref{eq:pgvwc07_ti_canonical_representation_sum}) is the scalar
    $\lambda^2$ of that block.
\end{proof}

\begin{lemma}[Uniqueness for a single block]
    \label{lem:pgvwc07_ti_uniqueness_single_block}
    \lean{MPSTensor.pgvwc07_uniqueness_of_unital,
        MPSTensor.pgvwc07_uniqueness_of_weighted,
        MPSTensor.sum_smul_mul_conjTranspose_smul,
        MPSTensor.sum_inv_smul_mul_conjTranspose_eq_one}
    \leanok
    \uses{def:l_blk_injective, def:mpv}
    Let $B$ and $C$ be tensors of bond dimension $D$ with
    $\sum_iB_iB_i^\dagger=\lambda_B^2\,\Id$ and
    $\sum_iC_iC_i^\dagger=\lambda_C^2\,\Id$ for some $\lambda_B,\lambda_C>0$,
    let $B$ be $L_0$-block injective, let $N>2L_0+D^4$, and let $B$ and $C$
    represent the same nonzero state on the ring of $N$ sites. Then
    $\lambda_B=\lambda_C$, and there are a unitary $U$ and a scalar
    $\xi\in\C$ with
    \begin{align}
        |\xi| & =1, & \xi^N & =1, & B_i & =\xi\,UC_iU^\dagger
        \qquad\text{for every }i.
        \label{eq:pgvwc07_ti_uniqueness_conclusion}
    \end{align}
\end{lemma}

\begin{proof}
    \leanok
    \uses{thm:pgvwc07_fixed_length_intertwiner,
        lem:pgvwc07_intertwiner_unitary}
    For $\lambda_B=\lambda_C=1$,
    Theorem~\ref{thm:pgvwc07_fixed_length_intertwiner} gives $R\neq0$ and
    $x\neq0$ with $RC_i=\tfrac1xB_iR$, and
    Lemma~\ref{lem:pgvwc07_intertwiner_unitary} turns $R$ into a unitary $U$
    with $|x|=1$ and $B_i=xUC_iU^\dagger$. The coefficients then satisfy
    $V^{(N)}(B)=x^NV^{(N)}(C)=x^NV^{(N)}(B)$, and since the state is nonzero,
    $x^N=1$; take $\xi=x$. For general weights, $\lambda_B^{-1}B$ and
    $\lambda_B^{-1}C$ represent the same state and the first is unital, so
    Theorem~\ref{thm:pgvwc07_fixed_length_intertwiner} gives $R$ and $x$ for
    this pair. The same $R$ satisfies
    $R\,(\lambda_C^{-1}C_i)=y^{-1}(\lambda_B^{-1}B_i)R$ with
    $y=x\lambda_C/\lambda_B$, and
    Lemma~\ref{lem:pgvwc07_intertwiner_unitary} applied to the two unital
    families gives $|y|=1$ and $B_i=xUC_iU^\dagger$. As before $x^N=1$, so
    $|x|=1$, and $|y|=1$ forces $\lambda_B=\lambda_C$.
\end{proof}

\begin{theorem}[Uniqueness of the canonical form, as printed]
    \label{thm:pgvwc07_ti_uniqueness_printed}
    \notready
    \uses{def:pgvwc07_ti_canonical_representation, def:l_blk_injective}
    Let $\lvert\psi\rangle=\sum\operatorname{tr}(B_{i_1}\cdots B_{i_N})
        \lvert i_1\cdots i_N\rangle$ be a translation-invariant canonical
    $D$-MPS such that condition C1 holds at some length $L_0$, the
    open-boundary canonical representation of $\lvert\psi\rangle$ is unique,
    and $N>2L_0+D^4$. Then, if $\lvert\psi\rangle$ admits another
    translation-invariant canonical $D$-MPS representation with matrices
    $C_i$, there is a unitary $U$ with $B_i=UC_iU^\dagger$ for every $i$.
    This is \cite[Theorem~7]{PerezGarcia2007Matrix} as printed. Its
    conclusion is false: for $d=D=1$, $B_0=1$ and $C_0=\eta$ with $\eta\neq1$
    an $N$th root of unity are two canonical representations of the same
    nonzero state on $N>3$ sites, and unitary conjugation acts trivially in
    bond dimension one \cite{gap:pgvwc07_quadratic_reconstruction_phase}.
    The corrected statement is Theorem~\ref{thm:pgvwc07_ti_uniqueness}.
\end{theorem}

\begin{theorem}[Uniqueness of the canonical form, corrected]
    \label{thm:pgvwc07_ti_uniqueness}
    \lean{MPSTensor.pgvwc07_uniqueness_of_canonical_form,
        MPSTensor.pgvwc07_uniqueness_of_canonical_form_rephased,
        MPSTensor.pgvwc07_uniqueness_of_canonical_form_unitGaugePhaseEquiv}
    \leanok
    \uses{def:pgvwc07_ti_canonical_representation, def:l_blk_injective,
        def:mpv}
    Let $B=(B_i)_i$ and $C=(C_i)_i$ be translation-invariant canonical
    representations of bond dimension $D$ in the sense of
    Definition~\ref{def:pgvwc07_ti_canonical_representation}. Assume that
    $B$ is $L_0$-block injective, that $N>2L_0+D^4$, and that $B$ and $C$
    represent the same nonzero state on the ring of $N$ sites. Then there
    are a unitary $U$ and a scalar $\xi\in\C$ satisfying
    (\ref{eq:pgvwc07_ti_uniqueness_conclusion}). The rephased
    representation $\widehat C_i:=\xi C_i$ represents the same length-$N$
    state and satisfies $B_i=U\widehat C_iU^\dagger$ exactly; equivalently,
    $C$ and $B$ are gauge-phase equivalent with the unitary gauge $U$ and the
    unit-modulus phase $\xi$.

    The printed hypothesis that the open-boundary canonical representation
    of the state is unique is not needed: the intertwiners it supplies in
    the printed proof are obtained from condition C1 alone in
    Theorem~\ref{thm:pgvwc07_fixed_length_intertwiner}. The corrected
    conclusion is the one recorded in \cite{gap:pgvwc07_ti_uniqueness_scope}.
\end{theorem}

\begin{proof}
    \leanok
    \uses{lem:pgvwc07_ti_canonical_representation_single_block,
        thm:pgvwc07_fixed_length_intertwiner,
        lem:pgvwc07_intertwiner_unitary,
        lem:pgvwc07_ti_uniqueness_single_block}
    The state is nonzero, so $D\ge1$, and
    Lemma~\ref{lem:pgvwc07_ti_canonical_representation_single_block} leaves
    $B$ with a single block: $\sum_iB_iB_i^\dagger=\lambda_B^2\,\Id$ with
    $0<\lambda_B\le1$. The tensors $\lambda_B^{-1}B$ and $\lambda_B^{-1}C$
    represent the same state and the first is unital and $L_0$-block
    injective, so Theorem~\ref{thm:pgvwc07_fixed_length_intertwiner} gives
    $R\neq0$ and $x\neq0$ with $R\,(\lambda_B^{-1}C_i)=\tfrac1x(\lambda_B^{-1}B_i)R$,
    and $R$ is invertible by Lemma~\ref{lem:pgvwc07_intertwiner_unitary}.
    Hence $C_i=\tfrac1xR^{-1}B_iR$ is conjugate to a nonzero multiple of
    $B_i$, so $C$ is $L_0$-block injective, and
    Lemma~\ref{lem:pgvwc07_ti_canonical_representation_single_block} leaves
    $C$ with a single block as well:
    $\sum_iC_iC_i^\dagger=\lambda_C^2\,\Id$ with $0<\lambda_C\le1$.
    Lemma~\ref{lem:pgvwc07_ti_uniqueness_single_block} now gives $U$ and
    $\xi$. Rephasing $C$ by $\xi$ multiplies every length-$N$ coefficient by
    $\xi^N=1$, and $\xi\overline{\xi}=|\xi|^2=1$ makes $(U,\xi)$ a
    unit-modulus gauge-phase equivalence.
\end{proof}

%% ---------------------------------------------------------
\section{One-sided inverse and physical realisation of virtual insertions}%
\label{sec:decomposition_map}
%% ---------------------------------------------------------

If $A$ is injective, the linear combination map
$\Phi_A(c):=\sum_i c_iA^i\colon\C^d\to\MN{D}$
is surjective and admits a right inverse. The resulting decomposition map
expresses any matrix as a linear combination of the spanning family
$\{A^i\}$, and its multiplicative refinement re-expresses a virtual insertion
on the bond as a physical operation on the local index.

\begin{lemma}[Existence of a right inverse]\label{lem:exists_rightInverse}
    \lean{Kraus.IsInjective.exists_rightInverse}
    \leanok
    \uses{def:injective}
    If $A$ is injective, then there exists a $\C$-linear map
    $\Psi\colon\MN{D}\to\C^d$ satisfying
    $\Phi_A\circ\Psi=\id_{\MN{D}}$.
\end{lemma}

\begin{proof}
    \leanok
    Surjectivity of $\Phi_A$ yields a linear section by choice.
\end{proof}

\begin{definition}[Decomposition map]\label{def:decomposition_map}
    \lean{Kraus.decompositionMap}
    \leanok
    \uses{def:injective, lem:exists_rightInverse}
    Let $A$ be an injective MPS tensor.
    The \emph{decomposition map}
    $\Psi_A\colon\MN{D}\to\C^d$ is a fixed choice of right inverse of
    $\Phi_A$.
\end{definition}

\begin{lemma}[Decomposition property]\label{lem:exists_decomposition}
    \lean{Kraus.IsInjective.exists_decomposition}
    \leanok
    \uses{def:decomposition_map}
    If $A$ is injective, then for every $X\in\MN{D}$ there exist
    coefficients $c\in\C^d$ such that
    $X=\sum_{i=0}^{d-1}c_iA^i$.
\end{lemma}

\begin{proof}
    \leanok
    Take $c=\Psi_A(X)$. Then
    $\Phi_A(c)=\Phi_A(\Psi_A(X))=X$.
\end{proof}

\begin{remark}[Non-uniqueness of the decomposition]
    \label{rem:decomposition_nonunique}
    When $d>D^2$, the map $\Phi_A$ has a nontrivial kernel, so
    the section $\Psi_A$ is not unique.
    The statements that follow depend only on
    $\Phi_A\circ\Psi_A=\id_{\MN{D}}$, so the choice does not
    affect the conclusions.
\end{remark}

\subsection{Physical realisation of virtual insertions}%
\label{sec:phys_realize}

For an injective tensor $A$, the physical realisation map encodes the
effect of inserting a virtual matrix $X$ on the bond as a $d\times d$
operation on the physical indices. Its multiplicativity
(Theorem~\ref{thm:physRealize_mul}) turns the bond algebra into an algebra
acting on the physical leg, and is the source of the algebra homomorphism
used in the virtual bond gauge.

\begin{definition}[Physical realisation map]\label{def:phys_realize}
    \lean{MPSTensor.physRealize}
    \leanok
    \uses{def:decomposition_map}
    Let $A$ be an injective tensor.
    For $X\in\MN{D}$, define the $d\times d$ matrix
    $O_A(X)\in\MN{d}$ by
    \begin{align}
        O_A(X)_{ij}
         & :=(\Psi_A(A^iX))_j,
        \notag
    \end{align}
    so that row~$i$ of $O_A(X)$ contains the decomposition coefficients
    of $A^iX$ in the spanning set $\{A^j\}$.
    Diagrammatically, the inserted virtual matrix is re-expressed as a
    physical operation on the same local tensor:
    \begin{center}
        % A^i X = sum_j O_A(X)_{ij} A^j: on the left, X is inserted on
        % A's east virtual leg.  On the right, the lower leg of O_A(X)
        % contracts with A's physical leg over j, while O_A(X)'s upper
        % leg is the free index i.  The operator carries no virtual
        % index.  Both sides have signature (1,1,1,0).
        \tenkzkernel
        \begin{tenkz}[cols=2, west=open, east=open]
            \tn[ports={90:physical:$i$}]{A} & \tn[skin=ring]{X}
        \end{tenkz}
        $ = $
        \begin{tenkz}[rows={op,ket}, cols=1]
            \tn[skin=mpo, ports={90:physical:$i$, 270:physical:$j$}]{O_A(X)} \\
            \tn[ports={90:physical, 180:virtual, 0:virtual}]{A}
        \end{tenkz}
    \end{center}
    The map $X\mapsto O_A(X)$ is $\C$-linear.
\end{definition}

\begin{lemma}[Defining property of the physical realisation map]%
    \label{lem:physRealize_spec}
    \lean{MPSTensor.physRealize_spec}
    \leanok
    \uses{def:phys_realize, def:decomposition_map}
    For every $X\in\MN{D}$ and every physical index $i$,
    \begin{align}
        A^iX
         & =\sum_{j=0}^{d-1}O_A(X)_{ij}A^j.
        \notag
    \end{align}
\end{lemma}

\begin{proof}
    \leanok
    By definition, $O_A(X)_{ij}=(\Psi_A(A^iX))_j$.
    Applying $\Phi_A$ gives
    $\sum_jO_A(X)_{ij}A^j=\Phi_A(\Psi_A(A^iX))=A^iX$,
    where the last step uses
    $\Phi_A\circ\Psi_A=\id_{\MN{D}}$.
\end{proof}

\begin{theorem}[Physical realisation is multiplicative]%
    \label{thm:physRealize_mul}
    \lean{MPSTensor.physRealize_mul}
    \leanok
    \uses{def:phys_realize, lem:physRealize_spec, def:injective}
    Let $A$ be an injective MPS tensor.
    For all $X,Y\in\MN{D}$,
    $O_A(XY)=O_A(X)O_A(Y)$.
\end{theorem}

\begin{proof}
    \leanok
    Applying Lemma~\ref{lem:physRealize_spec} to $X$ yields
    $A^iX=\sum_{j=0}^{d-1}O_A(X)_{ij}A^j$.
    Multiplying by $Y$ on the right and applying the linear map $\Psi_A$
    gives
    \begin{align}
        \Psi_A(A^iXY)
         & =\sum_{j=0}^{d-1}O_A(X)_{ij}\Psi_A(A^jY).
        \notag
    \end{align}
    Taking the $k$th component of both sides gives
    \begin{align}
        O_A(XY)_{ik}
         & =\sum_{j=0}^{d-1}O_A(X)_{ij}O_A(Y)_{jk}
        =(O_A(X)O_A(Y))_{ik}.
        \notag
    \end{align}
    Since this holds for all $i$ and $k$, the matrices are equal.
\end{proof}

\begin{remark}[General injective case]
    \label{rem:physRealize_mul_general}
    \uses{thm:physRealize_mul, def:injective, def:decomposition_map}
    When $d>D^2$, the decomposition map $\Psi_A$ need not be unique.
    Nevertheless, Theorem~\ref{thm:physRealize_mul} is already the
    general injective statement: for the fixed choice of $\Psi_A$ in
    Definition~\ref{def:decomposition_map}, the associated physical
    realisation map $O_A$ is multiplicative.
    This supersedes the earlier basis-only formulation.
\end{remark}

%% ---------------------------------------------------------
\section{The virtual bond gauge and the proportionality lemma}%
\label{sec:virtual_bond_gauge}
%% ---------------------------------------------------------

The two structural lemmas of~\cite{Molnar2018NormalPEPS} carry the entire
fixed-length argument. The virtual bond gauge
(Lemma~\ref{lem:virtual_bond_gauge}) turns equality of the combined MPV
families into the virtual-bond conjugacy $X\mapsto Z^{-1}XZ$ by a fixed
invertible matrix $Z$; the proportionality
lemma (Lemma~\ref{lem:tensor_proportional}) shows that two tensor pairs
agreeing under all virtual insertions are proportional. Both rest on two
elementary facts: injectivity is preserved by blocking, and the centre of
$\MN{D}$ consists of scalars.

\subsection{Two preliminary facts}
\label{sec:intermediate}

\begin{lemma}[Blocking preserves injectivity]\label{lem:blocking_preserves_injectivity}
    \lean{MPSTensor.chainCombinedTensor_isInjective}
    \leanok
    \uses{def:injective, def:blocked_tensor}
    Let $A_0,\ldots,A_{m-1}$ be injective MPS tensors with compatible
    bond dimensions.
    Then the blocked tensor
    \begin{align}
        (A_0\otimes\cdots\otimes A_{m-1})^{(i_0,\ldots,i_{m-1})}
         & :=A_0^{i_0}A_1^{i_1}\cdots A_{m-1}^{i_{m-1}}
        \notag
    \end{align}
    is again injective.
\end{lemma}

\begin{proof}
    \leanok
    The span of the blocked matrices contains all products
    $A_0^{i_0}\cdots A_{m-1}^{i_{m-1}}$.
    Since each $\{A_k^i\}_i$ spans $\MN{D}$, these products span
    all of $\MN{D}$.
\end{proof}

\subsection{Same state forces a virtual bond gauge}

The 3-site virtual-bond statement is used later in the proof of the PEPS
fundamental theorem. Its hypothesis is the all-length MPV equality of the
combined tensors, and its conclusion is the conjugacy of the middle-bond
insertion data.

\begin{lemma}[Virtual bond gauge -- 3-site case]%
    \label{lem:virtual_bond_gauge}
    \lean{MPSTensor.virtual_bond_gauge}
    \leanok
    \uses{thm:physRealize_mul, thm:simplicity}
    Let $A=(A_0,A_1,A_2)$ and $B=(B_0,B_1,B_2)$ be two
    3-site injective periodic chains with common bond dimension
    $D\ge1$ and $\mc{V}(A)=\mc{V}(B)$ (same MPV family).
    Then at the middle bond there exists an invertible matrix
    $Z\in\GL_D(\C)$ such that, for all $X\in\MN{D}$ and all physical
    indices $i_0,i_1,i_2$, the virtual-insertion trace coefficients agree
    after conjugation:
    \begin{align}
        \tr\!\left(A_0^{i_0}X A_1^{i_1}A_2^{i_2}\right)
         & =\tr\!\left(B_0^{i_0}Z^{-1}XZ B_1^{i_1}B_2^{i_2}\right).
        \notag
    \end{align}

    In \cite[Lemma~1]{Molnar2018NormalPEPS}, the statement is given for
    arbitrary chain length $n\ge3$ and any bond position, with
    uniqueness of $Z$ up to scalar.
    The chain-level generalization to arbitrary $n\ge3$
    and any bond position is not pursued here.
\end{lemma}

\begin{proof}
    \leanok
    \uses{thm:physRealize_mul, thm:simplicity, lem:blocking_preserves_injectivity}
    The $\mc{V}(A)=\mc{V}(B)$ hypothesis gives, for every $X\in\MN{D}$,
    equality of the virtual-insertion trace coefficients
    \begin{align}
        \tr\!\left(A_0^{i_0}X A_1^{i_1}A_2^{i_2}\right)
         & =\tr\!\left(B_0^{i_0}Y B_1^{i_1}B_2^{i_2}\right)
        \notag
    \end{align}
    for some $Y$ depending on $X$.
    Since all three site tensors on each side are injective,
    the physical realisation maps $O_A,O_B$ are algebra homomorphisms
    by Theorem~\ref{thm:physRealize_mul}.
    Setting $T(X)=Y$ defines a $\C$-linear multiplicative map,
    hence a nonzero $\C$-algebra endomorphism of $\MN{D}$, which is
    bijective by Theorem~\ref{thm:simplicity}.
    Skolem--Noether (Theorem~\cite[``Skolem--Noether'']{QICLean2026Blueprint}) gives
    $Z\in\GL_D(\C)$ with $T(X)=Z^{-1}XZ$.
\end{proof}

\subsection{Aligned gauges force proportionality}%
\label{sec:tensor_proportional}

After applying Lemma~\ref{lem:virtual_bond_gauge} to each bond and absorbing
the gauge matrices, one obtains two chains whose virtual-insertion
operations agree on every bond. Lemma~2 of~\cite{Molnar2018NormalPEPS}
shows that this agreement forces the local tensors to be proportional.

\begin{lemma}[Aligned gauges force proportionality --
        Lemma~2 of \protect\cite{Molnar2018NormalPEPS}]%
    \label{lem:tensor_proportional}
    \lean{MPSTensor.tensor_proportional}
    \leanok
    \uses{def:injective, def:decomposition_map, lem:center_scalar}
    Let $A_1,A_2$ and $B_1,B_2$ be injective MPS tensors with
    the same positive bond dimension $D$.
    Suppose that, for all $X\in\MN{D}$ on the internal bond,
    all $Y\in\MN{D}$ on the external bond, and all physical indices $i,j$,
    \begin{align}
        \tr\!\left(A_1^iX A_2^j\right)
         & =\tr\!\left(B_1^iX B_2^j\right),
        \label{eq:internal_bond}            \\
        \tr\!\left(A_2^jY A_1^i\right)
         & =\tr\!\left(B_2^jY B_1^i\right).
        \label{eq:external_bond}
    \end{align}
    Diagrammatically, inserting $X$ on the bond between the two sites,
    as in (\ref{eq:internal_bond}), gives
    \begin{center}
        % tr(A_1^i X A_2^j) = tr(B_1^i X B_2^j): X lies on the
        % internal bond between the two site tensors, and the periodic
        % return closes the remaining virtual ends.  The two upper legs
        % are the free indices i,j.  Both sides have signature (0,0,2,0).
        \tenkzkernel
        \begin{tenkz}[boundary=periodic]
            \tn[ports={90:physical:$i$}]{A_1} & \tn[skin=ring]{X} & \tn[ports={90:physical:$j$}]{A_2}
        \end{tenkz}
        $ = $
        \begin{tenkz}[boundary=periodic]
            \tn[ports={90:physical:$i$}]{B_1} & \tn[skin=ring]{X} & \tn[ports={90:physical:$j$}]{B_2}
        \end{tenkz}
    \end{center}
    while the insertion $Y$ on the outer trace bond, as
    in (\ref{eq:external_bond}), is
    \begin{center}
        % tr(A_2^j Y A_1^i) = tr(B_2^j Y B_1^i): reading each periodic
        % word from its second site encounters A_2/B_2, then Y on the
        % outer return bond, then A_1/B_1.  Thus this is not the internal
        % cyclic order above.  The upper legs are the free indices i,j,
        % and both sides have signature (0,0,2,0).
        \tenkzkernel
        \begin{tenkz}[boundary=periodic]
            \tn[ports={90:physical:$i$}]{A_1} & \tn[ports={90:physical:$j$}]{A_2} & \tn[skin=ring]{Y}
        \end{tenkz}
        $ = $
        \begin{tenkz}[boundary=periodic]
            \tn[ports={90:physical:$i$}]{B_1} & \tn[ports={90:physical:$j$}]{B_2} & \tn[skin=ring]{Y}
        \end{tenkz}
    \end{center}
    so the two hypotheses distinguish the internal and external virtual
    loops of the two-site periodic chain.

    Then there exists a nonzero scalar $\lambda\in\C^{\times}$ such
    that, for all physical indices $i,j$,
    \begin{align}
        A_1^i
         & =\lambda B_1^i,
        \notag                 \\
        A_2^j
         & =\lambda^{-1}B_2^j.
        \notag
    \end{align}
\end{lemma}

\begin{proof}
    \leanok
    \uses{def:decomposition_map, lem:center_scalar}
    By cyclicity and non-degeneracy of the trace pairing,
    the conditions in (\ref{eq:internal_bond})
    and (\ref{eq:external_bond}) give, for all physical indices $i,j$,
    \begin{align}
        A_2^jA_1^i
         & =B_2^jB_1^i,
        \label{eq:aft_product_identities_a} \\
        A_1^iA_2^j
         & =B_1^iB_2^j.
        \label{eq:aft_product_identities_b}
    \end{align}
    Write $\Id_D=\sum_jc_jA_2^j$ using the decomposition map
    $\Psi_{A_2}$, and define $Z:=\sum_jc_jB_2^j\in\MN{D}$.
    This is the point where the two-site comparison collapses to a single
    boundary matrix:
    \begin{center}
        % A_1^i = Z B_1^i = B_1^i Z: the upper leg is the free physical
        % index i, the red capsule Z acts on the west or east virtual
        % index, and the untouched outer virtual indices remain open.
        % All three terms have boundary signature (1,1,1,0).
        \tenkzkernel
        \begin{tenkz}[cols=1, west=open, east=open]
            \tn[ports={90:physical:$i$}]{A_1}
        \end{tenkz}
        $ = $
        \begin{tenkz}[cols=2, west=open, east=open]
            \tn[skin=ring]{Z} & \tn[ports={90:physical:$i$}]{B_1}
        \end{tenkz}
        $ = $
        \begin{tenkz}[cols=2, west=open, east=open]
            \tn[ports={90:physical:$i$}]{B_1} & \tn[skin=ring]{Z}
        \end{tenkz}


        % Z B_1^i = (lambda Id_D) B_1^i: after centrality identifies Z
        % with lambda Id_D, the two red capsules act on the same west
        % virtual index of B_1^i.  Both terms keep the east virtual index
        % and physical index i open, hence have signature (1,1,1,0).
        \begin{tenkz}[cols=2, west=open, east=open]
            \tn[skin=ring]{Z} & \tn[ports={90:physical:$i$}]{B_1}
        \end{tenkz}
        $ = $
        \begin{tenkz}[cols=2, west=open, east=open]
            \tn[skin=ring]{\lambda \Id_D} & \tn[ports={90:physical:$i$}]{B_1}
        \end{tenkz}
    \end{center}
    The first product identity (\ref{eq:aft_product_identities_a}) gives
    $A_1^i=ZB_1^i$, and the second
    (\ref{eq:aft_product_identities_b}) gives $A_1^i=B_1^iZ$.
    Hence $Z$ commutes with every $B_1^i$; since $\{B_1^i\}$ spans
    $\MN{D}$, $Z$ is central.
    By Lemma~\ref{lem:center_scalar}, $Z=\lambda\Id_D$ for some
    $\lambda\in\C^{\times}$ (nonzero since the tensors are nonzero).
    Thus $A_1^i=\lambda B_1^i$, and substituting
    into (\ref{eq:aft_product_identities_a}) forces
    $A_2^j=\lambda^{-1}B_2^j$.
\end{proof}

%% ---------------------------------------------------------
